\chapter{Prokofiev, Sergej}
\label{cha:prokofiev}

% \href{}{List of works} on Wikipedia.

\noindent{Symphonies – two juvenile: Symphony (1902) and Symphony (1908)} \\

\noindent{Ala i Lolli (1915, abandoned and rewritten as the Scythian Suite)} \\

\noindent{Sinfonietta in A (original version), Op. 5 (1909)} \textit{Royal Scottish National Orchestra, Neeme J\"{a}rvi}\marginnote{high} \\
\noindent{Dreams, Op. 6 (1910)} \\
\noindent{Two Poems for Female Chorus and Orchestra, Op. 7 (1909–10)} \\
\noindent{Autumnal, Op. 8 (1910)} \textit{Royal Scottish National Orchestra, Neeme J\"{a}rvi}\marginnote{high} \\

% 10
\noindent{Piano Concerto No. 1 in D$\flat$, Op. 10 (1911–12)} \textit{Jean-Efflam Bavouzet, BBC PO, Gianandrea Noseda}\marginnote{high} \\
\noindent{Piano Concerto No. 2 in G minor, Op. 16 (1912–13, lost, re-written in 1923)} \textit{Jean-Efflam Bavouzet, BBC PO, Gianandrea Noseda}\marginnote{high} \\
\noindent{The Ugly Duckling, Op. 18 (1914)} \\
\noindent{Violin Concerto No. 1 in D, Op. 19 (1916–17)} \textit{Janine Jansen, Oslo PO, Klaus M\"{a}kel\"{a}}\marginnote{max}  \\

% 20
\noindent{Scythian Suite, Op. 20 (from Ala i Lolli) (1914–15)} \textit{Royal Scottish National Orchestra, Neeme J\"{a}rvi}\marginnote{high} \\
\noindent{Chout / Tale of the Jester, Op. 21 (1915, rev. 1920)} \textit{Royal Scottish National Orchestra, Neeme J\"{a}rvi}\marginnote{high} \\
\noindent{Suite from Chout, Op. 21bis (1920)} \\
\noindent{Symphony No. 1 in D, Classical, Op. 25 (1916–17)} \textit{G\"{u}rzenich-Orchester K\"{o}ln, Dmitri Kitayenko}\marginnote{high} \\
\noindent{Piano Concerto No. 3 in C, Op. 26 (1917–21)} \textit{Jean-Efflam Bavouzet, BBC PO, Gianandrea Noseda}\marginnote{high} \\
\noindent{Andante from Piano Sonata No. 4, arranged for orchestra, Op. 29bis (1934)} \textit{Royal Scottish National Orchestra, Neeme J\"{a}rvi}\marginnote{high} \\
\noindent{Trapeze (1924, abandoned and rewritten as the Quintet for Oboe, Clarinet, Violin, Viola, and Double Bass)} \\

% 30
\noindent{Seven, They Are Seven, Op. 30 (1917–18, rev. 1933)} \\
\noindent{Suite from The Love for Three Oranges, Op. 33bis (1919, rev. 1924)} \textit{Royal Scottish National Orchestra, Neeme J\"{a}rvi}\marginnote{high} \\
\noindent{Overture on Hebrew Themes, Op. 34bis (1934)} \\
\noindent{Melodie, Op. 35bis (1920)} \\
\noindent{Vocal Suite from The Fiery Angel, Op. 37bis (1923, incomplete)} \\

% 40
\noindent{Symphony No. 2 in D minor, Op. 40 (1924–25)} \textit{G\"{u}rzenich-Orchester K\"{o}ln, Dmitri Kitayenko}\marginnote{high} \\
\noindent{Le pas d'acier / The Steel Step, Op. 41 (1926)} \\
\noindent{Suite from Le pas d'acier, Op. 41bis (1926)} \textit{Royal Scottish National Orchestra, Neeme J\"{a}rvi}\marginnote{high} \\
\noindent{American Overture, Op. 42 (1926) for 17 instruments; American Overture, Op. 42bis (1928) for full orchestra} \\
\noindent{Divertissement, Op. 43 (1925–29)} \textit{Royal Scottish National Orchestra, Neeme J\"{a}rvi}\marginnote{high} \\
\noindent{Symphony No. 3 in C minor, Op. 44 (1928)} \textit{G\"{u}rzenich-Orchester K\"{o}ln, Dmitri Kitayenko}\marginnote{high} \\
\noindent{Le fils prodigue / The Prodigal Son, Op. 46 (1929)} \textit{Royal Scottish National Orchestra, Neeme J\"{a}rvi}\marginnote{high} \\
\noindent{Suite from The Prodigal Son, Op. 46bis (1929)} \\
\noindent{Symphony No. 4 in C (original version), Op. 47 (1929–30)} \textit{G\"{u}rzenich-Orchester K\"{o}ln, Dmitri Kitayenko}\marginnote{high} \\
\noindent{Sinfonietta in A (revised version of Op. 5), Op. 48 (1929)} \textit{Lahti Symphony Orchestra, Dima Slobodeniuk}\marginnote{02/2022, max} \\
\noindent{Suite from The Gambler ("Four Portraits and Denouement"), Op. 49 (1931)} \\

% 50
\noindent{Andante from String Quartet No. 1, arranged for string orchestra, Op. 50bis (1930)} \textit{Royal Scottish National Orchestra, Neeme J\"{a}rvi}\marginnote{high} \\
\noindent{On the Dnieper, Op. 51 (1931)} \\
\noindent{Suite from On the Dnieper, Op. 51bis (1933)} \\
\noindent{Piano Concerto No. 4 in B$\flat$, Op. 53 (1931), for left hand} \textit{Jean-Efflam Bavouzet, BBC PO, Gianandrea Noseda}\marginnote{high} \\
\noindent{Piano Concerto No. 5 in G, Op. 55 (1932)} \textit{Jean-Efflam Bavouzet, BBC PO, Gianandrea Noseda}\marginnote{high} \\
\noindent{Symphonic Song, Op. 57 (1933)} \textit{Royal Scottish National Orchestra, Neeme J\"{a}rvi}\marginnote{high} \\
\noindent{Cello Concerto in E minor, Op. 58 (1933–38)} \\

% 60
\noindent{Lieutenant Kijé, Op. 60 (1934), also arranged as an orchestral suite} \\
\noindent{Suite from Lieutenant Kijé, Op. 60 (1934)} \textit{Royal Scottish National Orchestra, Neeme J\"{a}rvi}\marginnote{high} \\
\noindent{Suite from Egyptian Nights, Op. 61 (1938)} \\
\noindent{Violin Concerto No. 2 in G minor, Op. 63 (1935)}  \textit{Lisa Batiashvili, Chamber Orchestra of Europe, Yannick N\'{e}zet-S\'{e}guin}\marginnote{max} \\
\noindent{Romeo and Juliet, Op. 64 (1935, classified as a Soviet choreodrama, or drambalet)} \textit{LSO, Valery Gergiev}\marginnote{high}  \\
\noindent{Suite 1 from Romeo and Juliet, Op. 64bis (1936)} \\
\noindent{Suite 2 from Romeo and Juliet, Op. 64ter (1936)} \\
\noindent{Peter and the Wolf, Op. 67 (1936), a children's story for narrator and orchestra} \\

% 70
\noindent{The Queen of Spades, Op. 70a (1936), after Pushkin} \\
\noindent{Boris Godunov, Op. 70bis (1936)} \\
\noindent{Eugene Onegin, Op. 71 (1936)} \\
\noindent{Russian Overture, Op. 72 (1936)} \textit{Philharmonia Orchestra, Neeme J\"{a}rvi}\marginnote{high} \\
\noindent{Cantata for the 20th Anniversary of the October Revolution, Op. 74 (1936–37), cantata for 2 choruses, orchestra, military band, accordion band, and percussion band} \\
\noindent{Songs of Our Days, Op. 76 (1937), for chorus and orchestra} \\
\noindent{Hamlet, Op. 77 (1937–38)} \\
\noindent{Alexander Nevsky, Op. 78 (1939), cantata for mezzo-soprano, chorus, and orchestra} \textit{Olga Borodina, Kirov Chorus St Petersburg, Mariinsky Orchestra, Valery Gergiev}\marginnote{high}  \\
\noindent{Alexander Nevsky, Op. 78 (1938), film directed by Sergei Eisenstein (also exists in the form of a cantata)} \\
\noindent{Lermontov (1941)} \\
\noindent{Kotovsky (1942)} \\
\noindent{Tonya (1942)} \\
\noindent{The Partisans in the Ukrainian Steppes (1942)} \\



% 80
\noindent{Suite from Semyon Kotko, Op. 81bis (1941)} \\
\noindent{Zdravitsa, Op. 85 (1939), cantata for chorus and orchestra (also known as "Hail to Stalin")} \\
\noindent{Cinderella, Op. 87 (1944)} \\
\noindent{Symphonic March, Op. 88 (1941)} \\

% 90
\noindent{The Year 1941, Op. 90 (1941)} \\
\noindent{War and Peace, Suite, Op. 91} \textit{Philharmonia Orchestra, Neeme J\"{a}rvi}\marginnote{high} \\  
\noindent{Ballad of an Unknown Boy, Op. 93 (1942–43), for soloists, chorus, and orchestra} \\

% 100
\noindent{Symphony No. 5 in B$\flat$, Op. 100 (1944)} \textit{G\"{u}rzenich-Orchester K\"{o}ln, Dmitri Kitayenko}\marginnote{high} \\
\noindent{Suite 3 from Romeo and Juliet, Op. 101 (1946)} \\
\noindent{Ode to the End of the War, Op. 105 (1945)} \\
\noindent{Suite 1 from Cinderella, Op. 107 (1946)} \\
\noindent{Suite 2 from Cinderella, Op. 108 (1946)} \\
\noindent{Suite 3 from Cinderella, Op. 109 (1946)} \\



% 110
\noindent{Waltz Suite, Op. 110 (1946)} \textit{Royal Scottish National Orchestra, Neeme J\"{a}rvi}\marginnote{high} \\
\noindent{Symphony No. 6 in E$\flat$ minor, Op. 111 (1945–47)} \textit{G\"{u}rzenich-Orchester K\"{o}ln, Dmitri Kitayenko}\marginnote{high} \\
\noindent{Symphony No. 4 in C (revised version), Op. 112 (1947)} \textit{G\"{u}rzenich-Orchester K\"{o}ln, Dmitri Kitayenko}\marginnote{high} \\
\noindent{Thirty Years, Op. 113 (1947)} \\
\noindent{Flourish, Mighty Land, Op. 114 (1947), cantata for chorus and orchestra} \\

% 110
\noindent{Ivan the Terrible, Op. 116 (1942–45), film directed by Sergei Eisenstein (also exists in various concert forms)} \\
\noindent{Tale of the Stone Flower, Op. 118 (1953)} \textit{Royal Scottish National Orchestra, Neeme J\"{a}rvi}\marginnote{high} \\

% 120
\noindent{Pushkin Waltzes, Op. 120 (1949)} \\
\noindent{Winter Bonfire, Op. 122 (1949–50), for boy's choir and small orchestra} \\
\noindent{Summer Night, suite from Betrothal in a Monastery, Op. 123 (1950)} \textit{Philharmonia Orchestra, Neeme J\"{a}rvi}\marginnote{high} \\
\noindent{On Guard for Peace, Op. 124 (1950), cantata for chorus and orchestra} \\
\noindent{Symphony-Concerto for Cello and Orchestra in E minor, Op. 125 (1950–52)} \textit{Christian Polt\'{e}ra, Lahti SO, Anja Bihlmaier}\marginnote{07/2024, max} \\
\noindent{Wedding Suite from The Tale of the Stone Flower, Op. 126 (1951)} \textit{Royal Scottish National Orchestra, Neeme J\"{a}rvi}\marginnote{high} \\
\noindent{Gypsy Fantasy from The Tale of the Stone Flower, Op. 127 (1951)} \\
\noindent{Urals Rhapsody from The Tale of the Stone Flower, Op. 128 (1951)} \\
\noindent{The Mistress of Copper Mountain from The Tale of the Stone Flower, Op. 129 (1951); unfinished)} \\


% 130
\noindent{The Meeting of the Volga and the Don, Op. 130 (1951)} \\
\noindent{Symphony No. 7 in C$\sharp$ minor, Op. 131 (1951–52)} \textit{G\"{u}rzenich-Orchester K\"{o}ln, Dmitri Kitayenko}\marginnote{high} \\
\noindent{Cello Concertino in G minor, Op. 132 (1953–unfinished)} \\
\noindent{Piano Concerto No. 6, Op. 134 (1953–unfinished)} \\
\noindent{Symphony No. 2 in D minor (revised version), Op. 136 (1953; unrealized)} \\






